
\documentclass{article}
\usepackage{colortbl}
\usepackage{makecell}
\usepackage{multirow}
\usepackage{supertabular}

\begin{document}

\newcounter{utterance}

\centering \large Interaction Transcript for game `wordle', experiment `high\_frequency\_words\_no\_clue\_no\_critic', episode 1 with glm{-}5.2.
\vspace{24pt}

{ \footnotesize  \setcounter{utterance}{1}
\setlength{\tabcolsep}{0pt}
\begin{supertabular}{c@{$\;$}|p{.15\linewidth}@{}p{.15\linewidth}p{.15\linewidth}p{.15\linewidth}p{.15\linewidth}p{.15\linewidth}}
    \# & \multicolumn{2}{c}{Player} && \multicolumn{2}{c}{Game Master} \\
    \hline

    \theutterance \stepcounter{utterance}  
    & & & \multicolumn{4}{p{0.6\linewidth}}{
        \cellcolor[rgb]{0.9,0.9,0.9}{
            \makecell[{{p{\linewidth}}}]{
                \texttt{\tiny{[P1$\langle$GM]}}
                \texttt{Sei un mago del linguaggio a cui piace indovinare parole usando le regole fornite.} \\
\\ 
\texttt{Benvenuto a Wordle! Hai sei tentativi per indovinare la parola obiettivo, una parola valida di cinque lettere minuscole. Per favore usa i tag "spiegazione:" e "tentativo:" per fornire una spiegazione concisa per ogni tentativo.} \\
\\ 
\texttt{Per esempio, se il tuo tentativo è "parte", la tua risposta dovrebbe essere} \\
\texttt{spiegazione: questa è una comune parola di cinque lettere, e sto iniziando il mio tentativo con questa parola.} \\
\texttt{tentativo: parte} \\
\\ 
\texttt{Dopo ogni tentativo, la tua risposta verrà convalidata e riceverai un feedback che indica quali lettere sono corrette (verde), quali lettere sono corrette ma nella posizione sbagliata (giallo) e quali lettere sono errate (rosso). Questo feedback può essere utile per determinare quali lettere includere o escludere nel tuo prossimo tentativo.} \\
\\ 
\texttt{Per esempio, il feedback per "parte" potrebbe essere:} \\
\texttt{feedback\_tentativo: p<giallo> a<giallo> r<verde> t<giallo> e<rosso>} \\
\\ 
\texttt{La spiegazione dovrebbe contenere dettagli su come il feedback\_tentativo viene usato per arrivare a un nuovo tentativo.} \\
\\ 
\texttt{Non generare alcun altro testo.} \\
\\ 
\texttt{Iniziamo con il tuo primo tentativo.} \\
            }
        }
    }
     \\ \\

    \theutterance \stepcounter{utterance}  
    & \multicolumn{4}{p{0.6\linewidth}}{
        \cellcolor[rgb]{0.561,0.737,0.561}{
            \makecell[{{p{\linewidth}}}]{
                \texttt{\tiny{[P1$\rangle$GM]}}
                \texttt{spiegazione: inizio con una parola comune di cinque lettere che contiene vocali e consonanti frequenti nella lingua italiana.} \\
\texttt{tentativo: parte} \\
            }
        }
    }
    & & \\ \\

    \theutterance \stepcounter{utterance}  
    & & & \multicolumn{2}{p{0.3\linewidth}}{
        \cellcolor[rgb]{0.9,0.9,0.9}{
            \makecell[{{p{\linewidth}}}]{
                \texttt{\tiny{[GM$|$GM]}}
                \texttt{attempts = 1} \\
\texttt{target\_word = gente} \\
\texttt{guess = parte} \\
\texttt{guess\_feedback = p<red> a<red> r<red> t<green> e<green>} \\
            }
        }
    }
    & & \\ \\

    \theutterance \stepcounter{utterance}  
    & & & \multicolumn{4}{p{0.6\linewidth}}{
        \cellcolor[rgb]{0.9,0.9,0.9}{
            \makecell[{{p{\linewidth}}}]{
                \texttt{\tiny{[P1$\langle$GM]}}
                \texttt{feedback\_tentativo: p<rosso> a<rosso> r<rosso> t<verde> e<verde>} \\
\\ 
\texttt{Fornisci la tua risposta solo in questo formato.} \\
\texttt{spiegazione: dettagli} \\
\texttt{tentativo: parola} \\
            }
        }
    }
     \\ \\

    \theutterance \stepcounter{utterance}  
    & \multicolumn{4}{p{0.6\linewidth}}{
        \cellcolor[rgb]{0.561,0.737,0.561}{
            \makecell[{{p{\linewidth}}}]{
                \texttt{\tiny{[P1$\rangle$GM]}}
                \texttt{spiegazione: Il feedback conferma che le lettere "t" e "e" sono corrette e posizionate correttamente alla fine della parola. Le lettere "p", "a", "r" sono errate. Cerco una parola di cinque lettere che termini con "te" ed eviti le lettere escluse, includendo nuove consonanti e vocali comuni.} \\
\texttt{tentativo: vette} \\
            }
        }
    }
    & & \\ \\

    \theutterance \stepcounter{utterance}  
    & & & \multicolumn{2}{p{0.3\linewidth}}{
        \cellcolor[rgb]{0.9,0.9,0.9}{
            \makecell[{{p{\linewidth}}}]{
                \texttt{\tiny{[GM$|$GM]}}
                \texttt{attempts = 2} \\
\texttt{target\_word = gente} \\
\texttt{guess = vette} \\
\texttt{guess\_feedback = v<red> e<green> t<red> t<green> e<green>} \\
            }
        }
    }
    & & \\ \\

    \theutterance \stepcounter{utterance}  
    & & & \multicolumn{4}{p{0.6\linewidth}}{
        \cellcolor[rgb]{0.9,0.9,0.9}{
            \makecell[{{p{\linewidth}}}]{
                \texttt{\tiny{[P1$\langle$GM]}}
                \texttt{feedback\_tentativo: v<rosso> e<verde> t<rosso> t<verde> e<verde>} \\
\\ 
\texttt{Fornisci la tua risposta solo in questo formato.} \\
\texttt{spiegazione: dettagli} \\
\texttt{tentativo: parola} \\
            }
        }
    }
     \\ \\

    \theutterance \stepcounter{utterance}  
    & \multicolumn{4}{p{0.6\linewidth}}{
        \cellcolor[rgb]{0.561,0.737,0.561}{
            \makecell[{{p{\linewidth}}}]{
                \texttt{\tiny{[P1$\rangle$GM]}}
                \texttt{spiegazione: Il feedback indica che la parola termina con "tte". La prima 't' del tentativo precedente è risultata rossa, il che significa che la terza lettera non è una 't'. La 'e' è verde, indicando che la parola ha due 'e'. La struttura è "\_\_ette". Escludendo le lettere già provate, provo "nette".} \\
\texttt{tentativo: nette} \\
            }
        }
    }
    & & \\ \\

    \theutterance \stepcounter{utterance}  
    & & & \multicolumn{2}{p{0.3\linewidth}}{
        \cellcolor[rgb]{0.9,0.9,0.9}{
            \makecell[{{p{\linewidth}}}]{
                \texttt{\tiny{[GM$|$GM]}}
                \texttt{attempts = 3} \\
\texttt{target\_word = gente} \\
\texttt{guess = nette} \\
\texttt{guess\_feedback = n<yellow> e<green> t<red> t<green> e<green>} \\
            }
        }
    }
    & & \\ \\

    \theutterance \stepcounter{utterance}  
    & & & \multicolumn{4}{p{0.6\linewidth}}{
        \cellcolor[rgb]{0.9,0.9,0.9}{
            \makecell[{{p{\linewidth}}}]{
                \texttt{\tiny{[P1$\langle$GM]}}
                \texttt{feedback\_tentativo: n<giallo> e<verde> t<rosso> t<verde> e<verde>} \\
\\ 
\texttt{Fornisci la tua risposta solo in questo formato.} \\
\texttt{spiegazione: dettagli} \\
\texttt{tentativo: parola} \\
            }
        }
    }
     \\ \\

    \theutterance \stepcounter{utterance}  
    & \multicolumn{4}{p{0.6\linewidth}}{
        \cellcolor[rgb]{0.561,0.737,0.561}{
            \makecell[{{p{\linewidth}}}]{
                \texttt{\tiny{[P1$\rangle$GM]}}
                \texttt{spiegazione: Il feedback indica che la parola termina con "tte". La lettera 'n' è gialla, quindi appartiene alla parola ma non in prima posizione. La lettera 't' è rossa, il che è contraddittorio con il feedback precedente; si conferma che la terza lettera non è 't'. La struttura è "\_\_ette". Provo "bocche".} \\
\texttt{tentativo: bocche} \\
            }
        }
    }
    & & \\ \\

    \theutterance \stepcounter{utterance}  
    & & & \multicolumn{2}{p{0.3\linewidth}}{
        \cellcolor[rgb]{0.9,0.9,0.9}{
            \makecell[{{p{\linewidth}}}]{
                \texttt{\tiny{[GM$|$GM]}}
                \texttt{Error: The length of the guessed word is not 5.} \\
            }
        }
    }
    & & \\ \\

    \theutterance \stepcounter{utterance}  
    & & & \multicolumn{2}{p{0.3\linewidth}}{
        \cellcolor[rgb]{0.9,0.9,0.9}{
            \makecell[{{p{\linewidth}}}]{
                \texttt{\tiny{[GM$|$GM]}}
                \texttt{game\_result = ABORT} \\
            }
        }
    }
    & & \\ \\

\end{supertabular}
}

\end{document}
